\subsection{Связь опорных векторов и крайних точек}

\begin{theorem}[о совпадении опорного вектора и крайней точки]\label{th:support-extreme}
Пусть $x[N]$ — допустимая точка канонической задачи. Тогда
\[
  x[N] \ \text{— опорный вектор} \iff
  x[N] \ \text{— крайняя точка допустимого множества } S .
\]
\end{theorem}

\begin{idea}
В тетради в обоих пунктах написано $\Np$, но речь идёт именно о множестве
\emph{положительных} компонент: опорность по определению~\ref{def:support-vector}
проверяется для столбцов $\{a_j\}_{j \in \Nzp}$. Ниже везде пишем $\Nzp$ и $\Nzz$ —
это важно для корректности оценки $\theta > 0$ в первом пункте.
\end{idea}

\begin{proof}
\textbf{1) $(\Longleftarrow)$: крайняя точка $\implies$ опорный вектор.}
От противного: пусть $x[N]$ — \emph{не} опорный вектор. Тогда столбцы
$A[M,\Nzp]$ линейно зависимы, то есть
\[
  \exists\, u[\Nzp] \ne 0 : \qquad A[M, \Nzp]\, u[\Nzp] = 0 .
\]
Доопределим $u$ нулями вне $\Nzp$, то есть положим $u[i] = 0$ для всех
$i \in \Nzz$; тогда $A[M,N]\,u[N] = 0$.

Положим
\[
  x_1[N] = x[N] - \theta\, u[N], \qquad
  x_2[N] = x[N] + \theta\, u[N],
\]
где $\theta = \min\{\theta_1, \theta_2\}$ и
\[
  \theta_1 = \min_{\substack{i \in \Nzp \\ u[i] > 0}} \frac{x[i]}{u[i]} ,
  \qquad
  \theta_2 = \min_{\substack{i \in \Nzp \\ u[i] < 0}} \frac{-x[i]}{u[i]} .
\]
(Если один из наборов индексов пуст, соответствующий минимум считаем равным
$+\infty$.)

По подбору $\theta$ векторы $x_1$ и $x_2$ удовлетворяют ограничению на знак:
на $\Nzz$ обе поправки нулевые, а на $\Nzp$ сдвиг не выводит за ноль ни в одну,
ни в другую сторону. Кроме того $\theta > 0$, так как $x[i] > 0$ при $i \in \Nzp$,
и нулевых числителей там нет.

Проверим допустимость:
\begin{align*}
  A[M,N]\,x_1[N] &= A[M,N]\,x[N] - \theta\,\underbrace{A[M,N]\,u[N]}_{=\,0} , \\
  A[M,N]\,x_2[N] &= A[M,N]\,x[N] + \theta\,\underbrace{A[M,N]\,u[N]}_{=\,0} ,
\end{align*}
откуда
\[
  A[M,N]\,x_1[N] = A[M,N]\,x[N] = b[M] = A[M,N]\,x_2[N] .
\]
Значит $x_1, x_2 \in S$, причём $x_1 \ne x_2$ (так как $\theta > 0$ и $u \ne 0$), и
\[
  x[N] = \tfrac{1}{2}\big(x_1[N] + x_2[N]\big) .
\]
Получили противоречие с тем, что $x$ — крайняя точка.

\medskip
\textbf{2) $(\Longrightarrow)$: опорный вектор $\implies$ крайняя точка.}
От противного: пусть $x[N]$ — \emph{не} крайняя точка. Тогда
\[
  \exists\, x_1[N], x_2[N] \in S,\ x_1 \ne x_2 : \qquad
  x[N] = \tfrac{1}{2}\big(x_1[N] + x_2[N]\big) .
\]
Так как $x_1, x_2 \ge 0$, а на множестве $\Nzz$ полусумма равна нулю, обе точки
обращаются там в ноль:
\[
  x[\Nzz] = 0 \implies x_1[\Nzz] = 0, \quad x_2[\Nzz] = 0 .
\]
Поэтому ограничения для них сводятся к системам только по базисным индексам:
\[
  \begin{cases}
    A[M, \Nzp]\, x_1[\Nzp] = b[M], \\
    A[M, \Nzp]\, x_2[\Nzp] = b[M] .
  \end{cases}
\]
Вычитая, получаем
\[
  A[M, \Nzp]\big(x_1[\Nzp] - x_2[\Nzp]\big) = 0 ,
\]
причём $x_1[\Nzp] \ne x_2[\Nzp]$ (вне $\Nzp$ обе точки нулевые, так что различаться
они могут только там). Значит столбцы $A[M,\Nzp]$ линейно зависимы, то есть
$x[N]$ — не опорный вектор. Противоречие.
\end{proof}

\inputfig{figures/fig-theta-segment}

\subsection{Оптимальное решение достигается в опорном векторе}

\begin{theorem}\label{th:optimal-is-support}
Пусть $\bar{x}_*[N]$ — оптимальное решение канонической задачи \emph{с минимальным
числом положительных компонент}. Тогда $\bar{x}_*[N]$ — опорный вектор.
\end{theorem}

\begin{proof}
Рассмотрим разбиение индексов на $\Nzp$ и $\Nzz$ для точки $\bar{x}_*$. От
противного: пусть $\bar{x}_*$ — не опорный вектор. Тогда, как в
теореме~\ref{th:support-extreme},
\[
  \exists\, u[\Nzp] \ne 0 : \quad A[M,\Nzp]\, u[\Nzp] = 0 ,
\]
доопределённое нулями до $u[N]$, так что $A[M,N]\,u[N] = 0$.

Положим
\[
  x_\theta[N] = \bar{x}_*[N] - \theta\, u[N] .
\]
Эта точка допустима при всех $\theta \in [-\theta_0, \theta_0]$, где
\[
  \theta_0 = \min_{\substack{i \in \Nzp \\ u[i] \ne 0}}
             \frac{\bar{x}_*[i]}{\abs{u[i]}} > 0 :
\]
равенства $A x_\theta = b$ выполнены автоматически (поправка лежит в ядре $A$),
а ограничение на знак — по выбору $\theta_0$.

Значение целевой функции:
\[
  c\T[N]\, x_\theta[N] = c\T[N]\, \bar{x}_*[N] - \theta\, c\T[N]\, u[N] .
\]
Так как $x_\theta$ допустима при \emph{обоих} знаках $\theta$, знак $\theta$ можно
выбрать в зависимости от знака $c\T[N]\,u[N]$ так, чтобы
$\theta\, c\T[N]\,u[N] \ge 0$, то есть чтобы целевая функция не возрастала. Но по
определению оптимального решения не может быть
\[
  c\T[N]\, x_\theta[N] < c\T[N]\, \bar{x}_*[N] ,
\]
следовательно
\[
  c\T[N]\, u[N] = 0 \quad\implies\quad
  c\T[N]\, x_\theta[N] = c\T[N]\, \bar{x}_*[N]
  \quad \forall \theta \in [-\theta_0, \theta_0] ,
\]
то есть $x_\theta$ — тоже оптимальное решение.

Возьмём теперь $\theta = \theta_0$. Минимум в определении $\theta_0$ достигается на
некотором индексе $i_0 \in \Nzp$, и для него $x_{\theta_0}[i_0] = 0$. Остальные
компоненты остались неотрицательными, а компоненты из $\Nzz$ — нулевыми. Значит
$x_{\theta_0}$ — оптимальное решение, у которого положительных компонент
\emph{строго меньше}, чем у $\bar{x}_*$. Противоречие с минимальностью.
\end{proof}

\begin{corollary*}
Оптимальное решение можно найти, перебрав все крайние точки (= опорные векторы) и
посчитав в каждой значение целевой функции. По лемме~\ref{lem:count} перебор конечен.
\end{corollary*}

Но такой перебор может быть вычислительно трудным: число вариантов $\binom{n}{m}$
растёт экспоненциально.
